\corrigesChapitre{6}{Divisibilité et congruences}
\begin{corrige}{1}{Diviseurs}
\begin{questions}
\item $1,2,3,6,9,18$.
\item Oui : $18=(-3)(-6)$.
\item Non : $0k=0$ pour tout entier $k$. En revanche tout entier divise zéro.
\end{questions}
\end{corrige}
\begin{corrige}{2}{Division euclidienne}
\begin{questions}
\item $83=7\times11+6$.
\item $-83=7\times(-12)+1$.
\item Les restes $6$ et $1$ sont dans $\{0,\ldots,6\}$.
\end{questions}
\end{corrige}
\begin{corrige}{3}{Lire une congruence}
\begin{questions}
\item $29-1=28$ est divisible par $7$.
\item $4$, car $-3=7(-1)+4$.
\item $2+5k$, avec $k\in\Z$.
\end{questions}
\end{corrige}
\begin{corrige}{4}{Calcul de restes}
\begin{questions}
\item $24\equiv0$, $35\equiv5$, donc reste $5$.
\item $0\times5\equiv0$.
\item $25\equiv1\pmod6$.
\end{questions}
\end{corrige}
\begin{corrige}{5}{Combinaison entière}
\begin{questions}
\item Un diviseur commun divise toute combinaison entière.
\item La combinaison vaut $1$, donc $d=1$.
\item Ils sont premiers entre eux, quel que soit l’entier $n$.
\end{questions}
\end{corrige}
\begin{corrige}{6}{Puissance modulo $7$}
\begin{questions}
\item $2,4,1$.
\item $100=6\times16+4$, donc reste $4$.
\item La fonction rend $4$ ; le calcul modulaire explique le résultat.
\end{questions}
\end{corrige}
\begin{corrige}{7}{Critère modulo $9$}
\begin{questions}
\item $10^k\equiv1\pmod9$, donc le nombre et la somme de ses chiffres ont même reste modulo $9$.
\item $7+3+2+6=18$, donc le nombre est divisible par $9$.
\item Il établit une congruence, pas une égalité au chiffre zéro. Un nombre non nul peut avoir reste zéro.
\end{questions}
\end{corrige}
\begin{corrige}{8}{Critère modulo $11$}
\begin{questions}
\item La somme alternée des chiffres, depuis les unités, doit être divisible par $11$.
\item $1-3+3-1=0$ : divisible par $11$.
\item $4-3+2-1=2$ : non divisible par $11$.
\end{questions}
\end{corrige}
\begin{corrige}{9}{Ne pas diviser sans condition}
\begin{questions}
\item $8-2=6$.
\item Cela donnerait $1\equiv4\pmod6$, faux puisque $6$ ne divise pas $3$.
\item $6\mid2(1-4)$ implique $3\mid1-4$, donc $1\equiv4\pmod3$.
\end{questions}
\end{corrige}
\begin{corrige}{10}{Reste d’un polynôme}
\begin{questions}
\item $n^3-n=n(n-1)(n+1)$ contient deux entiers consécutifs, donc un facteur pair.
\item Parmi trois entiers consécutifs, l’un est multiple de $3$.
\item Le produit est divisible par $2$ et $3$, qui sont premiers entre eux ; donc par $6$ (ou repérer directement les facteurs $2$ et $3$ dans le produit).
\end{questions}
\end{corrige}
\begin{corrige}{11}{Un cycle décimal}
\begin{questions}
\item $7^2\equiv9$, donc $7^4\equiv81\equiv1$.
\item $2026\equiv2\pmod4$, donc $7^{2026}\equiv7^2\equiv9\pmod{10}$.
\item Elle réduit l’exposant grâce au cycle de longueur $4$, et non parce qu’on pourrait réduire arbitrairement un exposant modulo $10$.
\end{questions}
\end{corrige}
\begin{corrige}{12}{Deux nombres consécutifs}
\begin{questions}
\item Le diviseur commun divise leur différence $(n+1)-n=1$.
\item Le PGCD positif est $1$.
\item Il divise alors $2$, donc il peut valoir $1$ ou $2$ ; la parité de $n$ les distingue.
\end{questions}
\end{corrige}
